% !PRINTED 405
\weberpage{405}
%<<CONT
\gesperrt{von der Art, dass $u, u_1, \ldots, u_{\nu-1}$ rational durch $\vartheta$
allein ausgedrückt werden können.}

Um ihn zu beweisen, setzen wir

\weq{1}{u = \dfrac{\varphi(t)}{\psi(t)}, \quad u_1 = \dfrac{\varphi_1(t)}{\psi_1(t)}, \quad \cdots, \quad u_{\nu-1} = \dfrac{\varphi_{\nu-1}(t)}{\psi_{\nu-1}(t)},}

und verstehen unter $\varphi(t)$, $\psi(t)$ ganze Functionen von $t$
ohne gemeinschaftlichen Theiler; ebenso unter $\varphi_1(t)$,
$\psi_1(t)$ u.\,s.\,f. Ausserdem dürfen wir noch voraussetzen, dass
von den Functionen $u, u_1, \ldots, u_{\nu-1}$ keine eine Constante sei.
Wir führen neben der Variablen $t$ eine zweite Variable $\tau$ ein und
setzen:

\weq{2}{v = \dfrac{\varphi(\tau)}{\psi(\tau)}, \quad v_1 = \dfrac{\varphi_1(\tau)}{\psi_1(\tau)}, \quad \cdots, \quad v_{\nu-1} = \dfrac{\varphi_{\nu-1}(\tau)}{\psi_{\nu-1}(\tau)}.}

Nun bilden wir die ganzen Functionen der beiden Variablen $t, \tau$,
deren keine identisch verschwindet:

\weq{3}{\begin{aligned}
\varphi(t)\,\psi(\tau) - \varphi(\tau)\,\psi(t) &= \Phi(t,\tau) = -\Phi(\tau,t)\\
\varphi_1(t)\,\psi_1(\tau) - \varphi_1(\tau)\,\psi_1(t) &= \Phi_1(t,\tau) = -\Phi_1(\tau,t)\\
&\cdots\cdots\cdots\cdots\cdots\cdots\cdots\cdots\cdots\cdots\cdots\cdots\cdots\cdots\\
\varphi_{\nu-1}(t)\,\psi_{\nu-1}(\tau) - \varphi_{\nu-1}(\tau)\,\psi_{\nu-1}(t) &= \Phi_{\nu-1}(t,\tau) = -\Phi_{\nu-1}(\tau,t).
\end{aligned}}

Betrachten wir sie als Functionen von $t$, so verschwinden sie alle,
wenn $t = \tau$ wird, und sie haben also einen grössten
gemeinschaftlichen Theiler, der in Bezug auf $t$ mindestens vom
ersten Grade ist, und vom $\mu^{\text{ten}}$ Grade sein möge:

\weq{4}{R = R(t, \tau).}

Diese Function $R$ ist auch in Bezug auf $\tau$ rational, und wenn
man sie so einrichtet, dass $\tau$ nicht im Nenner und nicht in einem
überflüssigen Factor vorkommt, so kann $R(t, \tau)$ bei der
Vertauschung von $t$ und $\tau$ höchstens sein Vorzeichen ändern.
Denn die Functionen (3) sind, auch als Functionen der beiden
Variablen $t, \tau$ betrachtet, durch $R(t, \tau)$, und da sie
alternirend sind, durch $R(\tau, t)$ theilbar (Bd.\,I, \S.\,51). Es
ist also $R(t, \tau)$ durch $R(\tau, t)$ theilbar, und umgekehrt, und
folglich unterscheiden sich beide nur durch einen constanten Factor,
der $= \pm 1$ sein muss, weil die nochmalige Vertauschung von $t$ mit
$\tau$ die ursprüngliche Function wieder herstellt.

Es lässt sich noch beweisen, dass keine der Functionen (3) (bei
unbestimmtem $\tau$) als Function von $t$ betrachtet, einen Factor
%CONT>>
